\newpage
\subsection{\BreakoutName}

Two clusterings are examined for \BreakoutName, obtained at layer 4 with $K=2$ and at layer 6b with $K=5$.
At the shallower layer, with only two clusters, the partition reduces to a coarse temporal split between the intact wall at the beginning of the episode and the partially broken wall observed once bricks have been cleared.
At the deeper layer, with five clusters, the partition is more fine-grained: it separates the beginning of the episode into three clusters distinguished mainly by the paddle's position, and isolates two more advanced game states, one in which the agent has dug a tunnel through the left side of the wall and one in which the ball is behind the wall.
The coarser clustering achieves a higher \gls{acu} than the finer one, at the cost of the more specific semantics it identifies.

\begin{table}[H]
    \centering
    \begin{tblr}{
        colspec={|c|c|c|c|c|},
        hlines,
        vlines,
        row{1} = {font=\bfseries, halign=c, valign=m}
    }
        Layer & $K$ & \gls{cu}       & \gls{as}       & \gls{acu}      \\
        4     & 2   & 0.99 $\pm$0.00 & 0.88 $\pm$0.00 & 0.87 $\pm$0.00 \\
        6b    & 5   & 0.79 $\pm$0.03 & 0.82 $\pm$0.01 & 0.61 $\pm$0.01 \\
    \end{tblr}
    \caption{\gls{cu}, \gls{as}, and \gls{acu} for the two clusterings analyzed in the \BreakoutName environment.}
    \label{tab:qualitative_breakout_metrics}
\end{table}

\begin{figure}[H]
    \qualitativeClusterFigureSetup
    \qualitativeClusterTwo
        {figures/clustering_agent_representation/pictures/breakout/qualitative_alt/cluster_0/image_1.png}
        {figures/clustering_agent_representation/pictures/breakout/qualitative_alt/cluster_0/image_2.png}
        {A}
    \qualitativeClusterTwo
        {figures/clustering_agent_representation/pictures/breakout/qualitative_alt/cluster_1/image_1.png}
        {figures/clustering_agent_representation/pictures/breakout/qualitative_alt/cluster_1/image_2.png}
        {B}
    \caption{Clusters A--B identified by \gls{car} in the \BreakoutName environment (layer 4, $K=2$).}
    \label{fig:qualitative_breakout_alt}
\end{figure}


\begin{figure}[H]
    \qualitativeClusterFigureSetup
    \qualitativeClusterTwo
        {figures/clustering_agent_representation/pictures/breakout/qualitative/cluster_0/image_1.png}
        {figures/clustering_agent_representation/pictures/breakout/qualitative/cluster_0/image_2.png}
        {C}
    \qualitativeClusterTwo
        {figures/clustering_agent_representation/pictures/breakout/qualitative/cluster_1/image_1.png}
        {figures/clustering_agent_representation/pictures/breakout/qualitative/cluster_1/image_2.png}
        {D}
    \qualitativeClusterTwo
        {figures/clustering_agent_representation/pictures/breakout/qualitative/cluster_2/image_1.png}
        {figures/clustering_agent_representation/pictures/breakout/qualitative/cluster_2/image_2.png}
        {E}
    \qualitativeClusterTwo
        {figures/clustering_agent_representation/pictures/breakout/qualitative/cluster_3/image_1.png}
        {figures/clustering_agent_representation/pictures/breakout/qualitative/cluster_3/image_2.png}
        {F}
    \qualitativeClusterTwo
        {figures/clustering_agent_representation/pictures/breakout/qualitative/cluster_4/image_1.png}
        {figures/clustering_agent_representation/pictures/breakout/qualitative/cluster_4/image_2.png}
        {G}
    \caption{Clusters C--G identified by \gls{car} in the \BreakoutName environment (layer 6b, $K=5$).}
    \label{fig:qualitative_breakout}
\end{figure}

